The project also contains a cross-domain specialist-routing extension.
Its question is different from the main paper's prefix-action question:
given two fixed within-domain specialists, can a shared learned router choose between them better than a hard domain split?

The answer is yes in expectation once supervision is large enough.
In the current extension results, learned routing can surpass hard specialist routing in mean regret.
At the same time, routing-family winners remain much less stable than the within-domain controller rankings used in the main text.
That is why routing stays in the appendix.
It is informative, but it lives at a different level of control.

Three conclusions matter for the paper.
First, hard-coded fallback heuristics are not the main source of routing gains.
Second, learned routing can outperform hard specialist routing in mean performance.
Third, routing must be reported with mean/std and win counts rather than with a single-seed winner, because single-run frontier effects are too unstable.

The current strongest routing family is a higher-capacity random-forest-style router, specifically the stronger ExtraTrees-style variant within that family.
We keep this result as an extension because it supports a broader hierarchical-control view, but it is not necessary for the benchmark and controller-class claims in the main paper.
